\section{Related work}
\label{sec:related}

Telperion stands on Lean~4 and its mathematical library, Mathlib
\cite{lean4,mathlib}: the emitted proofs are Lean, and the kernel that checks them
is Lean's. The idea of producing a checkable certificate rather than trusting a
procedure is the reflection/certificate tradition in proof assistants, and in the
positivity setting specifically it goes back to Positivstellensatz and
sum-of-squares certification \cite{positivstellensatz,lasserre_sos}, which several
of Telperion's shapes emit.

Two recent efforts from Axiom Math are close in spirit, and we want to keep them
distinct because they are distinct projects. \emph{AXLE} \cite{axle} is a cloud
Lean-verification utility; we drew engineering patterns (verification, gap-filling,
repair, negative controls) from it, but no code. \emph{AxiomMath/ZetaZeros}
is a Lean formalization of a pair-correlation zero-simplicity result of Lamzouri
\cite{axiommath_zetazeros}; two of Telperion's emitters port a \emph{proof idea}
from that formalization---a Montgomery--Taylor extremum-at-the-boundary move and a
transcendental enclosure atom---independently re-implemented in our idiom. Conflating the
verification utility with the formalization would misattribute both; the
repository's \texttt{NOTICE} file records the split in full.

The framing of this paper stands on a classical lineage we want to cite
plainly. The explicit formula connecting zeros to primes descends from
Riemann's memoir through Guinand \cite{guinand1948} and Weil \cite{weil1952};
the quasicrystal reading of that identity---the zeros as a one-dimensional
point set with pure-point diffraction on the prime powers, and the suggestion
that classifying such sets is a road toward RH---is Dyson's \cite{dyson2009}.
Our contribution to that lineage is deliberately modest and precisely scoped:
the \emph{finite-height} form of the identity, and the supporting localization
and region results, as kernel-checked certificates. The classical analytic work
of estimating boundary terms and passing to limits is not re-done here, and the
classification program itself we leave exactly where Dyson left it.

For independent checking, we use the Comparator \cite{comparator} from OpenAI's
ten-proofs project \cite{tenproofs}, together with \texttt{nanoda} \cite{nanoda},
a from-scratch Lean type checker in Rust; agreement between independently written
checkers is the basis of our strongest confidence claim. In the results and excursion sections, the
mathematics is due to others and cited there: Brualdi and Goldwasser for the tree
problem \cite{brualdi1984}; de la Vallée-Poussin for the classical zero-free region
\cite{dlvp1896}; Grigoriev and Laurent, with the symmetric-SOS technique of
Kurpisz, Lepp\"anen and Mastrolilli, for the proof-complexity lower bound
\cite{grigoriev2001,klm2016}.
